%%B000001 p

C++是最受欢迎、使用最广泛的编程语言之一，三十年来一直如此。C++的设计着眼于性能、效率与灵活性，融合了面向对象、命令式、泛型以及晚近的函数式编程等多种范式。C++由国际标准化组织（ISO）标准化，过去十年间经历了巨大的变化。随着C++11的标准化，这门语言进入了一个新时代，人们普遍称之为现代C++。类型推导、移动语义、lambda表达式、智能指针、统一初始化、可变参数模板等众多新特性改变了我们编写C++代码的方式，以至于它看上去几乎像是一门全新的编程语言。


%%B000002 p

本书涵盖C++11、C++14以及即将到来的C++17中的许多新特性。全书按案例组织，每个案例讨论一项特定的语言或库特性，或者开发者常遇到的一个问题及其使用现代C++的典型解法。通过100多个案例，你将学会掌握语言核心特性与标准库，包括字符串、容器、算法、迭代器、输入输出、正则表达式、线程、文件系统、原子操作与各类工具。


%%B000003 p

本书的写作历时约6个月，在此期间C++17标准的工作也在推进。在撰写本前言时，该标准已经完成，但尚未获批，预计今年晚些时候发布。书中有若干案例涵盖了C++17特性，包括折叠表达式、constexpr if、结构化绑定、新的标准属性，以及optional、any、variant与string\_view类型，还有filesystem库。


%%B000004 p

书中的所有案例都包含代码示例，展示如何使用某项特性或如何解决某个问题。这些代码示例是用Visual Studio 2017编写的，但也用Clang和GCC编译过。由于各编译器对各种语言与库特性的支持是逐步加入的，建议你使用最新版本，以确保所有这些特性都得到支持。撰写本书时，最新的版本是GCC 7.0、Clang 5.0和VC++ 2017（版本19.1）。尽管GCC与Clang支持本书涉及的全部特性，VC++尚不支持折叠表达式、constexpr if，以及std::search() 的搜索器（searchers）。


%%B000005 h3

\subsection*{本书内容}


%%B000006 p

\texttt{第1章}，\textit{现代C++语言核心特性}介绍现代C++语言核心特性，包括类型推导、统一初始化、限定作用域枚举、基于范围的for循环、结构化绑定等。


%%B000007 p

\texttt{第2章}，\textit{数字与字符串}讨论数值与字符串之间的转换、生成伪随机数、使用正则表达式，以及各种字符串类型。


%%B000008 p

\texttt{第3章}，\textit{函数探究}深入讲解默认化与删除的函数、可变参数模板、lambda表达式与高阶函数。


%%B000009 p

\texttt{第4章}，\textit{预处理与编译}考察编译的各个方面，从如何做条件编译，到编译期断言、代码生成，以及用属性向编译器提供提示。


%%B000010 p

\texttt{第5章}，\textit{标准库容器、算法与迭代器}向你介绍若干标准容器与大量算法，并教你编写自己的随机访问迭代器。


%%B000011 p

\texttt{第6章}，\textit{通用工具}深入讲解chrono库；any、optional与variant类型；并带你了解类型特征。


%%B000012 p

\texttt{第7章}，\textit{文件与流}说明如何读写流中的数据、用I/O操纵器控制流，并探索filesystem库。


%%B000013 p

\texttt{第8章}，\textit{线程与并发}告诉你如何使用线程、互斥量、锁、条件变量、promise、future与原子类型。


%%B000014 p

\texttt{第9章}，\textit{健壮性与性能}聚焦异常、常量正确性、类型转换、智能指针与移动语义。


%%B000015 p

\texttt{第10章}，\textit{模式与惯用法}涵盖多种有用的模式与惯用法，例如pimpl惯用法、非虚接口惯用法，以及奇异递归模板模式。


%%B000016 p

\texttt{第11章}，\textit{测试框架}帮你快速上手三种最广泛使用的测试框架：Boost.Test、Google Test与Catch。


%%B000017 h3

\subsection*{阅读本书所需}


%%B000018 p

本书中的代码可以从你在https://\allowbreak{}www.\allowbreak{}packtpub.\allowbreak{}com/ 的账户下载，不过我建议你自己动手把全部示例写一遍。要编译这些代码，在Windows上需要VC++ 2017，在Linux与Mac上需要GCC 7.0或Clang 5.0。如果你没有最新版本的编译器，或者想试试别的编译器，可以使用一个在线可用的编译器。虽然可用的在线平台有很多，但我推荐在GCC与Clang上使用https://\allowbreak{}wandbox.\allowbreak{}org/\allowbreak{}，在VC++上使用http://\allowbreak{}webcompiler.\allowbreak{}cloudapp.\allowbreak{}net/\allowbreak{}。


%%B000019 h3

\subsection*{本书面向的读者}


%%B000020 p

本书面向所有C++开发者，无论其经验水平如何。典型读者是入门级或中级C++开发者，他们希望掌握这门语言，成为多产的现代C++开发者。有经验的C++开发者则可以把本书当作一份很好的参考，其中涉及的许多C++11、C++14与C++17语言与库特性时常能派上用场。本书包含一百多个案例，难度分为简单、中级与高级。不过，它们都要求你先具备C++基础知识，包括函数、类、模板、命名空间、宏等。因此，如果你对这门语言还不熟悉，建议先读一本入门书熟悉核心内容，然后再来读本书。


%%B000021 h2

\section*{小节说明}


%%B000022 p

在本书中，你会看到几个反复出现的标题（准备工作、具体做法、工作原理、更多内容与延伸阅读）。


%%B000023 p

为了清楚地说明如何完成一个案例，我们按如下方式使用这些小节：


%%B000024 h3

\subsection*{准备工作}


%%B000025 p

本小节告诉你案例中会涉及什么，并说明如何准备案例所需的软件或预先设置。


%%B000026 h3

\subsection*{具体做法}


%%B000027 p

本小节包含按该案例操作所需的步骤。


%%B000028 h3

\subsection*{工作原理}


%%B000029 p

本小节通常详细解释上一小节中发生的事情。


%%B000030 h3

\subsection*{更多内容}


%%B000031 p

本小节包含关于该案例的补充信息，让读者对案例有更多了解。


%%B000032 h3

\subsection*{延伸阅读}


%%B000033 p

本小节提供指向该案例其他有用信息的有用链接。


%%B000034 h3

\subsection*{本书约定}


%%B000035 p

在本书中，你会看到多种文本样式，用以区分不同类型的信息。下面是这些样式的示例及其含义说明。


%%B000036 p

正文中的代码词、文件夹名、文件名、文件扩展名、路径名、URL、用户输入等按如下方式显示：\myttbr{\#include <\allowbreak{}iostream>}。


%%B000037 p

代码块按如下方式排版：


%%B000038 code

\begin{codeblockbg}
\begin{Verbatim}[commandchars=\\\{\},breaklines=true,breakanywhere=true,fontsize=\codefont,xleftmargin=17bp]
\PYG{c+cp}{\PYGZsh{}}\PYG{c+cp}{include}\PYG{+w}{ }\PYG{c+cpf}{\PYGZlt{}iostream\PYGZgt{}}
\PYG{k+kt}{int}\PYG{+w}{ }\PYG{n+nf}{main}\PYG{p}{(}\PYG{p}{)}
\PYG{p}{\PYGZob{}}
\PYG{+w}{  }\PYG{n}{std}\PYG{o}{:}\PYG{o}{:}\PYG{n}{cout}\PYG{+w}{ }\PYG{o}{\PYGZlt{}}\PYG{o}{\PYGZlt{}}\PYG{+w}{ }\PYG{l+s}{\PYGZdq{}}\PYG{l+s}{Hello World!}\PYG{l+s}{\PYGZdq{}}\PYG{+w}{ }\PYG{o}{\PYGZlt{}}\PYG{o}{\PYGZlt{}}\PYG{+w}{ }\PYG{n}{std}\PYG{o}{:}\PYG{o}{:}\PYG{n}{endl}\PYG{p}{;}
\PYG{+w}{  }\PYG{k}{return}\PYG{+w}{ }\PYG{l+m+mi}{0}\PYG{p}{;}
\PYG{p}{\PYGZcb{}}
\end{Verbatim}
\end{codeblockbg}


%%B000039 note

\begin{myNote}
警告或重要说明会出现在这样的框里。
\end{myNote}


%%B000040 note

\begin{myNote}[tip]
提示与技巧会这样显示。
\end{myNote}


%%B000041 h2

\section*{读者反馈}


%%B000042 p

我们始终欢迎读者的反馈。请告诉我们你对本书的看法——你喜欢什么，不喜欢什么。读者反馈对我们很重要，它能帮助我们开发出让你真正学有所获的图书。


%%B000043 p

如果要给我们发送一般性反馈，只需发邮件到\myttbr{feedback@packtpub.\allowbreak{}com}，并在邮件主题中注明书名。


%%B000044 p

如果你在某个领域有专长，并且有兴趣写书或为图书做贡献，请参阅我们的作者指南\myttbr{www.\allowbreak{}packtpub.\allowbreak{}com/\allowbreak{}authors}。


%%B000045 h2

\section*{客户支持}


%%B000046 p

现在你已经是Packt图书的用户了，我们准备了不少东西来帮助你充分利用所购图书。


%%B000047 h3

\subsection*{下载示例代码}


%%B000048 p

你可以从http://\allowbreak{}www.\allowbreak{}packtpub.\allowbreak{}com上的账户下载本书的示例代码文件。如果你是在别处购买的本书，可以访问http://\allowbreak{}www.\allowbreak{}packtpub.\allowbreak{}com/\allowbreak{}support注册，让这些文件直接邮寄给你。


%%B000049 p

你可以按以下步骤下载代码文件：


%%B000050 p

使用你的电子邮箱地址和密码登录或注册我们的网站。


%%B000051 p

将鼠标指针悬停在顶部的\textbf{SUPPORT}标签上。


%%B000052 p

点击\textbf{Code Downloads \& Errata}。


%%B000053 p

在\textbf{Search}框中输入书名。


%%B000054 p

选择你要下载代码文件的那本书。


%%B000055 p

在下拉菜单中选择你购买本书的渠道。


%%B000056 p

点击\textbf{Code Download}。


%%B000057 p

你也可以点击Packt Publishing网站上本书页面中的\textbf{Code Files}按钮来下载代码文件。在该页面的\textbf{Search}框中输入书名即可访问该页面。请注意，你需要登录自己的Packt账户。


%%B000058 p

文件下载完成后，请确保使用以下工具的最新版本解压该文件夹：


%%B000059 li

\begin{itemize}
\item Windows下的WinRAR / 7-Zip


%%B000060 li

\item Mac下的Zipeg / iZip / UnRarX


%%B000061 li

\item Linux下的7-Zip / PeaZip

\end{itemize}


%%B000062 p

本书的代码包也托管在GitHub上：https://\allowbreak{}github.\allowbreak{}com/\allowbreak{}PacktPublishing/\allowbreak{}Modern-\allowbreak{}Cpp-\allowbreak{}Programming-\allowbreak{}Cookbook。我们丰富的图书与视频目录中还有其他代码包，位于https://\allowbreak{}github.\allowbreak{}com/\allowbreak{}PacktPublishing/\allowbreak{}，去看看！


%%B000063 h3

\subsection*{勘误}


%%B000064 p

尽管我们已尽一切努力确保内容的准确性，但错误仍在所难免。如果你在我们某本书中发现了错误——也许是正文的错误，也许是代码的错误——我们将非常感谢你向我们报告。这样做可以让其他读者免于困扰，也有助于我们改进本书的后续版本。如果你发现勘误，请访问http://\allowbreak{}www.\allowbreak{}packtpub.\allowbreak{}com/\allowbreak{}submit-\allowbreak{}errata报告：选择你的书，点击Errata Submission Form链接，并填写勘误的详细信息。你的勘误一经核实，提交就会被接受，勘误会上传到我们的网站，或加入该书的Errata小节下已有的勘误列表。


%%B000065 p

要查看此前提交的勘误，请访问https://\allowbreak{}www.\allowbreak{}packtpub.\allowbreak{}com/\allowbreak{}books/\allowbreak{}content/\allowbreak{}support，在搜索框中输入书名。所需信息会显示在Errata小节下。


%%B000066 h3

\subsection*{盗版}


%%B000067 p

互联网上对受版权保护材料的盗版是各类媒体长期存在的问题。在Packt，我们非常重视对版权与许可证的保护。如果你在互联网上发现我们作品的任何形式的非法副本，请立即向我们提供其位置地址或网站名称，以便我们采取补救措施。


%%B000068 p

请通过\myttbr{copyright@packtpub.\allowbreak{}com}联系我们，并附上疑似盗版材料的链接。


%%B000069 p

感谢你帮助我们保护作者，以及我们为你带来有价值内容的能力。


%%B000070 h3

\subsection*{问题}


%%B000071 p

如果你对本书的任何方面有疑问，可以通过\myttbr{questions@packtpub.\allowbreak{}com}联系我们，我们会尽力解决。
